\subsection{Evaluable expressions and effective presentations}
\label{subsec:p0-evaluation-effectivity}

The early formulation of AEG began with rational arithmetic and explicitly
restricted its subsequent discussion to \emph{evaluable expressions}
\cite{Yuan2022AEGEvaluableExpressions}.  A later foundational question connected
this restriction with the difficulty of deciding whether a real expression
has a value \cite{Yuan2023AEGFoundationalProblem}.  The distinction matters:
defining an expression, defining its partial interpretation, computing its
value, and deciding its admissibility are different tasks.

\begin{definition}[Evaluability at an input]
\label{def:p0-evaluable-expression}
For a finite one-hole arithmetic expression $e[z]$ over $K$, let $D_e\subseteq K$
be the set of inputs at which every node admits ordinary arithmetic
evaluation, including every intermediate division guard.  Write
\[
  \nu_e:D_e\longrightarrow K
\]
for that evaluation.  The expression is \emph{evaluable at $x$} precisely when
$x\in D_e$.  A closed expression is \emph{evaluable} when its entire recursive
evaluation is defined.  These are semantic conditions on an already generated
syntax tree; they do not assert the existence of a decision procedure.
\end{definition}

For a comb, $D_e$ is exactly the final prefix domain defined above.  A
syntactically well-formed tree may have an empty domain: the expression
$z/0$ is one example.  This is partiality of evaluation, rather than failure
of the tree grammar or incompleteness of a proof system.  Alternatively, one
may require evidence of defined operands and a nonzero denominator in the
formation rule for division; such a rule defines a language of certified
evaluable expressions.  Checking supplied evidence and automatically finding
it remain separate tasks.

\begin{example}[A cancelled reciprocal retains its excluded input]
\label{ex:p0-cancelled-reciprocal-domain}
Let $e[z]=1/(1/z)$.  Its first reciprocal requires $z\ne0$; for such $z$, the
value $1/z$ is nonzero, so the second reciprocal is defined and returns $z$.
Thus
\[
  D_e=K\setminus\{0\},\qquad \nu_e(z)=z\quad(z\in D_e).
\]
The expression $z$ itself has domain $K$.  The two expressions agree on $D_e$
but do not define the same partial map.  Recording the pair $(D_e,\nu_e)$
preserves this difference after algebraic cancellation.  The missing input is
not a pole of the simplified identity function.  Composing the two projective
reciprocals instead permits an intermediate infinity and gives the identity
on $\mathbb P^1(K)$; that is a different interpretation, not an ordinary
evaluation at zero.  Neither exclusion by itself asserts a singularity of
the regular AES structure introduced in the next section.
\end{example}

The choice $K=\mathbb Q$ supplies a useful effective starting point when
constants and inputs are given as integer fractions with nonzero
denominators.  With unbounded exact integer arithmetic, evaluation of a finite
closed tree, or of a finite tree at a supplied rational input, terminates with
either a rational value or a division-by-zero failure: there are finitely many
nodes, and each arithmetic operation and zero test is effective.  This does
not give every syntax tree a value.  A bounded implementation must also
distinguish a resource refusal from arithmetic undefinedness.

The rational carrier is being used here for its effective presentation,
not for a discrete topology: $\mathbb Q$ is not discrete in the topology
inherited from $\mathbb R$.  The algebraic definitions of this section allow
an arbitrary field $K$ without an effective presentation; finite notation
cannot, however, name every real constant over a fixed finite alphabet.
The real geometric models impose their own smoothness hypotheses.  Passing
to such a model supplies neither exact numerical algorithms nor a general
test for admissibility.

Richardson's theorem concerns decision problems in a richer function
language.  With rational constants, $\pi$, $\log 2$, the variable, exponential
and sine, closed under addition, subtraction, multiplication and composition,
the existence of a negative value is undecidable; adding absolute value also
makes functional identity with zero undecidable
\cite{Richardson1968UndecidableProblems}.  This is not a prohibition on
defining real expressions, nor a theorem that every language of closed real
constants has undecidable zero tests.  The signature, presentation and
question must be specified before transferring such an obstruction.

Adva develops the finite construction viewpoint by retaining typed program
structure and construction history, and separating exact structural checks
from later numerical observations \cite{Yuan2026AdvaProgramProcessCore}.
The rational examples above explain one effective starting point for that
viewpoint.  Finite syntax and checkable evidence do not entail that every
semantic question is decidable, and a bounded checker does not replace the
continuous geometry.  Keeping domains and histories explicit allows these
discrete constructions and continuous interpretations to be studied together.
